% PROMOTED after independent root and lie_theory audits, 2026-11-19 18:21 UTC. Main part_B.tex is authoritative; this retains draft provenance.
\begin{saresult}{B.082}{Proved}{A universal coefficient-syzygy criterion for the two polynomial tensors}
Let $R\subset S$ be any unital inclusion of commutative rings. The
following construction uses only ideals, coefficient functions and
finite linear relations associated with this inclusion. It imposes
no finiteness, reducedness, flatness or domain hypothesis.

First form
\[
\begin{gathered}
 C=\bigcup_{d\ge1}(R:_S D_d(R)),\qquad \mathcal F=C[X],\\
 E=\operatorname{im}\left(\mathcal F\xrightarrow{\varepsilon}
                       \operatorname{App}(R,C/R)\right),\\
 \varepsilon\left(\sum c_iX^i\right)(r)=\sum c_ir^i+R.
\end{gathered}
                                                        \tag{B82.1}
\]
The first arithmetic test is that every $c\in C$ have a finite
matrix presentation
\[
\begin{gathered}
 c=Y M Z,\qquad M\in\operatorname{Mat}_{n}(R),\\
 Y\in C^{1\times n},\quad Z\in C^{n\times1},\quad
 YM\in R^{1\times n},\quad MZ\in R^{n\times1}.
\end{gathered}                                           \tag{B82.2}
\]
This test is equivalent to epicness of $R\to C$.
If it fails, the requested tensors are not uniquely determined.
When it holds, define the following finite-depth syzygy invariants.
Set
\[
\begin{aligned}
 U_q&=\operatorname{Tor}_q^R(\mathcal F,\mathcal F),\\
 V_q&=\operatorname{Tor}_q^R(\mathcal F,E)
                    \oplus\operatorname{Tor}_q^R(E,\mathcal F),\\
 W_q&=\operatorname{Tor}_q^R(E,E),\qquad
 U_q\xrightarrow{a_q}V_q\xrightarrow{b_q}W_q,
\end{aligned}                                           \tag{B82.3}
\]
where
\[
\begin{aligned}
 a_q&=(\operatorname{Tor}_q(1,\varepsilon),
                         \operatorname{Tor}_q(\varepsilon,1)),\\
 b_q&=\operatorname{Tor}_q(\varepsilon,1)
                        -\operatorname{Tor}_q(1,\varepsilon).
\end{aligned}
                                                        \tag{B82.4}
\]
The two terms of $b_q$ act on the respective summands of $V_q$;
thus $b_qa_q=0$. Put
\[
 Z_q=\ker a_q,\quad T_q=\operatorname{coker}b_q,\quad
 K_1=\ker b_1/\operatorname{im}a_1.                       \tag{B82.5}
\]
There are canonical maps, explicitly constructed in the proof,
\[
\begin{gathered}
 \delta_0:C[X,Y]=Z_0\longrightarrow T_1,
 \qquad \delta_1:Z_1\longrightarrow T_2,\\
 K_2=T_2/\operatorname{im}\delta_1.
\end{gathered}                                           \tag{B82.6}
\]
The two polynomial tensors on 153/004 exist uniquely if and only if
(B82.2) holds and
\[
\begin{gathered}
                         K_1=K_2=0,\\
 \delta_0\left(\sum_{i,j\ge0}\binom{i+j}{i}
                        a_{i+j}X^iY^j\right)=0,\\
 \delta_0\left(\sum_{i\ge0}a_iX^iY^i\right)=0            \tag{B82.7}
\end{gathered}
\]
for every finite coefficient vector $(a_i)$ in $C$ satisfying
$\sum_i a_ir^i\in R$ for every $r\in R$. Missing coefficients are zero.
All quantities in this test are defined before forming any tensor
power of $\Omega(R,S)$; in particular it is not a test that presupposes
the image or kernel of that tensor map. Its free-module implementation
below uses only finite coefficient relations in homological degrees
at most three. The criterion does not assert a uniformly terminating
algorithm for arbitrary presentations of infinite rings, and it
does not replace the separate higher-coherence conditions for a
plethory by a binary test.
\end{saresult}
\begin{proof}
We first check the matrix test and then construct the invariants and
prove their meaning by explicit chain calculations.
If (B82.2) holds, balancing its entries shows
$c\otimes1=1\otimes c$ in $C\otimes_R C$: use $YM$ to move the first
sum to the second factor, and then $MZ$ to move it back.
Conversely suppose this tensor equality holds. A finite presentation
of that equality gives finitely many $s_i\in C$, with labeled
$s_0=1,s_1=c$, and a finite matrix $(m_{ij})$ over $R$ such that
\[
 \sum_i m_{i0}s_i=c,\qquad \sum_i m_{i1}s_i=-1,
 \qquad \sum_i m_{ij}s_i=0\ (j\ne0,1),\qquad
                       \sum_jm_{ij}s_j=0\ (\text{all }i).
                                                        \tag{B82.8}
\]
Indeed present the second copy of $C$ by the free module on its
elements. The vector having entries $c$ at $s_0$, $-1$ at $s_1$,
and zero elsewhere lies in the image of its relation module tensored
with the first copy of $C$. Only finitely many relations and elements
are needed, giving (B82.8), after enlarging the common index set.
Remove the row and column indexed by zero, and call the remaining
matrix $M'$, with row $Y'=(s_i)_{i\ne0}$ and column
$Z'=(s_j)_{j\ne0}$. Then $Y'M'$ and $M'Z'$ have entries in $R$,
and $Y'M'Z'=m_{00}-c$. Appending a one-by-one block $m_{00}$ to
$-M'$, and a coordinate $1$ to both $Y'$ and $Z'$, gives (B82.2).
This proves the equivalence, including the finite certificate.
It is the opened Stacks Lemma 10.107.11 and Remark 10.107.12
\cite{StacksEpicMatrices};
the argument also supplies the final elementary step omitted there.
B.064 now proves that failure of this test precludes uniqueness.

Assume henceforth (B82.2). Since $C\otimes_R C\simeq C$, one has
$C\otimes_R(C/R)=0$. The module $E$ is a quotient of $(C/R)[X]$,
and $\mathcal F\to E$ is surjective. Right exactness therefore gives
\[
 \mathcal F\otimes_R E=E\otimes_R\mathcal F=0,
                         \qquad E\otimes_R E=0.          \tag{B82.9}
\]
In particular $V_0=W_0=0$, while
$U_0=\mathcal F\otimes_R\mathcal F=C[X,Y]$ by coefficient
multiplication. Hence $Z_0=U_0$ as claimed.

Here is a fixed, functorial free construction, so that the
transgressions in (B82.6) are unambiguous. For any $R$-module $M$,
let $P_0(M)$ be the free $R$-module on its underlying set, augmented
by $[m]\mapsto m$. Take $P_1(M)$ to be free on the elements of the
kernel, with its evident map to $P_0(M)$, and repeat on successive
kernels. A module map sends the indicated generators to their
images at each stage, making $P_\bullet$ functorial. Every generator
of a kernel is a finite linear relation in the preceding free
module. Thus this construction is explicit even without finite
generation or any choice of a tensor presentation of $\Omega$.

Write $P=P_\bullet(\mathcal F)$, $Q=P_\bullet(E)$, and
$u=P_\bullet(\varepsilon):P\to Q$. For nonnegative vertical degree
$q$, form three complexes of free $R$-modules
\[
\begin{aligned}
 \mathcal T_{0,\bullet}&=\operatorname{Tot}(P\otimes_RP),\\
 \mathcal T_{1,\bullet}&=\operatorname{Tot}(P\otimes_RQ)
                         \oplus\operatorname{Tot}(Q\otimes_RP),\\
 \mathcal T_{2,\bullet}&=\operatorname{Tot}(Q\otimes_RQ).
\end{aligned}                                           \tag{B82.10}
\]
Here $\operatorname{Tot}(P\otimes Q)_q=
\bigoplus_{a+b=q}P_a\otimes_RQ_b$, with vertical differential
$d=d_P\otimes1+(-1)^a1\otimes d_Q$ on that summand.
Define the horizontal chain maps
\[
 h_0=(1\otimes u,u\otimes1),\qquad
 h_1=(u\otimes1,-1\otimes u).                            \tag{B82.11}
\]
They commute with $d$ and satisfy $h_1h_0=0$.
Their induced maps on vertical homology are exactly $a_q,b_q$.
The vertical homology identifications with Tor follow by the free
double-resolution argument, as in opened Stacks Lemma 10.75.5
\cite{StacksTor}.

For $q=0,1$ and $[z]\in Z_q$, choose a vertical cycle
$z\in\mathcal T_{0,q}$ representing it. Since $a_q[z]=0$,
there exists $w\in\mathcal T_{1,q+1}$ such that
\[
                dw=(-1)^{q+1}h_0z.
\]
The element $(-1)^{q+1}h_1w$ is a vertical cycle in
$\mathcal T_{2,q+1}$. Define
\[
        \delta_q[z]=[(-1)^{q+1}h_1w]
                  \pmod{\operatorname{im}b_{q+1}}\in T_{q+1}.
                                                        \tag{B82.12}
\]
Changing $w$ by a vertical cycle changes this class by an image of
$b_{q+1}$. Changing $z$ to $z+dv$, with
$v\in\mathcal T_{0,q+1}$, permits replacing $w$ by
$w+(-1)^{q+1}h_0v$, which leaves $h_1w$ unchanged.
This proves independence of choices and $R$-linearity. Standard
comparison maps, or the same calculation with homotopic chain maps,
identify these definitions for any other free resolutions; the fixed
functorial resolution already defines them without that ambiguity.
For $q=0$ an arbitrary coefficient polynomial $H\in C[X,Y]$ is
represented by any preimage in $\mathcal T_{0,0}$, and (B82.12)
specifies its obstruction using one finite chain lift $w$.

We now identify what these invariants measure. Give
$\mathcal T_{p,q}$ total homological degree $q-p$, and set
\[
                      D=d+(-1)^q h.                     \tag{B82.13}
\]
The commutation $dh=hd$ and $h^2=0$ give $D^2=0$.
This total complex has homology in degree $n$ equal to
$\operatorname{Tor}_n^R(A,A)$, where $A=\ker\varepsilon=\Omega(R,S)$,
and zero homology in negative degrees. We verify this assertion
without a convergence assumption. The two-column total complex
$L$ formed from $P\xrightarrow{u}Q$, in degrees $q-p$ with
$p=0,1$, has $L_n=P_n\oplus Q_{n+1}$ and differential (B82.13).
The exact sequence $0\to A\to\mathcal F\to E\to0$ shows
$H_0(L)=A$ and all other homology zero. Also $L_n=0$ for $n<-1$,
and $L_{-1}=Q_0$ is free. The surjection $L_0\to L_{-1}$ therefore
splits, separating a contractible two-term complex. The remainder
is a nonnegative projective resolution of $A$: its degree-zero
module is a direct summand of the free module $L_0$, and its
higher modules are free. Tensoring two copies of this decomposition
shows that $L\otimes_RL$ computes $\operatorname{Tor}_\bullet^R(A,A)$.

There is an explicit sign isomorphism from $L\otimes_RL$ to the
three-column complex (B82.10)--(B82.13). On the component whose
first factor has column $p_1\in\{0,1\}$ and whose second factor
has vertical resolution degree $b$, multiply by $(-1)^{p_1b}$.
For the second vertical differential this changes its sign
$(-1)^{a-p_1}$ to $(-1)^a$; for the first horizontal differential
it changes $(-1)^a$ to $(-1)^{a+b}$; for the second horizontal
differential it leaves $(-1)^{a+b-p_1}$, which is the sign in
(B82.11)--(B82.13). The first vertical sign is unchanged.
This verifies the chain isomorphism term by term and proves the
homology assertion. In particular
\[
                 H_0(\operatorname{Tot}\mathcal T)=A\otimes_RA,
                 \qquad H_{-1}(\operatorname{Tot}\mathcal T)=0.
                                                        \tag{B82.14}
\]
The edge map obtained by taking the column-zero component and
then its vertical homology is the natural map
$A\otimes_RA\to\mathcal F\otimes_R\mathcal F=C[X,Y]$.
Indeed $H_0(L)=A$ is identified with the kernel of $\varepsilon$
by the augmentation to $\mathcal F$. Thus this edge map is
precisely $\theta_2$.

For completeness the low-degree exactness can now be checked
directly. A total degree-zero cycle consists of
\[
 x_0\in\mathcal T_{0,0},\quad
 x_1\in\mathcal T_{1,1},\quad
 x_2\in\mathcal T_{2,2},\qquad
             dx_1=-h_0x_0,\quad dx_2=h_1x_1.             \tag{B82.15}
\]
It follows from (B82.12) that its edge value lies in
$\ker\delta_0$. Conversely, if $\delta_0[x_0]=0$, choose $x_1$
solving the first equation. The zero class of $-h_1x_1$ modulo
$\operatorname{im}b_1$ allows adding a vertical cycle to $x_1$
so that $h_1x_1$ becomes a vertical boundary; then choose $x_2$
solving the second equation. Therefore
$\operatorname{im}\theta_2=\ker\delta_0$.
If $y\in\mathcal T_{2,1}$ is a vertical cycle, it is a total
degree-minus-one cycle. By (B82.14), it is the boundary of a
triple $(x_0,x_1,x_2)$ of total degree zero. The column-one
equation is $dx_1=-h_0x_0$, and the column-two equation is
$-h_1x_1+dx_2=y$. Thus $\delta_0[x_0]=[y]$ in $T_1$.
Every class in $T_1$ has such a representative, proving surjectivity
and the canonical cokernel identification
\[
                    \operatorname{coker}\theta_2\simeq T_1.
                                                        \tag{B82.16}
\]

If a total degree-zero cycle has zero edge value, its $x_0$ is a
vertical boundary. Subtracting a total boundary makes $x_0=0$.
Then $x_1$ is a vertical cycle, and (B82.15) says that its class
lies in $\ker b_1$. Different ways to remove $x_0$ change that
class by an image of $a_1$. This defines a canonical surjection
\[
                       \ker\theta_2\longrightarrow K_1. \tag{B82.17}
\]
Surjectivity follows by starting with a representative of a class
in $\ker b_1$ and lifting $h_1x_1$ to $dx_2$.
If the class in $K_1$ is zero, subtracting further total boundaries
makes $x_1=0$ as well. What remains is a vertical cycle
$x_2\in\mathcal T_{2,2}$. Such a pure column-two cycle is a total
boundary exactly when, modulo a vertical boundary, it is of the form
$h_1w$, where
\[
 dz=0,\quad z\in\mathcal T_{0,1},\qquad
 w\in\mathcal T_{1,2},\quad dw=h_0z.                     \tag{B82.18}
\]
Indeed the column-zero and column-one components of a boundary
of total degree one must vanish, giving precisely these equations;
the column-two component is $h_1w$ plus a vertical boundary.
Changing $w$ by a cycle contributes $\operatorname{im}b_2$,
and (B82.18) is exactly the defining lift for $\delta_1$.
Consequently (B82.17) has kernel $K_2$, with a canonical exact
sequence
\[
             0\longrightarrow K_2\longrightarrow\ker\theta_2
                     \longrightarrow K_1\longrightarrow0.   \tag{B82.19}
\]
This proves that uniqueness is equivalent to $K_1=K_2=0$.

All coefficient vectors of integer-valued polynomials belong to
$C$ by B.064. The two explicit arrays in (B82.7) are exactly the
coefficient polynomials of $F(X+Y)$ and $F(XY)$.
Equations (B82.15)--(B82.16) say that they lift exactly when their
$\delta_0$ classes vanish. If (B82.19) is zero, these lifts are
unique and recover the requested tensors. This proves both
directions of (B82.7) for every original inclusion $R\subset S$.

Every complex used to define the criterion is a free complex built
from $\mathcal F$, $E$ and $\varepsilon$, which themselves come
directly from (B82.1). Relations in the free resolutions and the
lifts in (B82.12), (B82.15), and (B82.18) have finite support.
Only $U_1,V_1,W_1$, $V_2,W_2$ and boundaries in vertical degree
three are needed. Thus the test is an explicit arithmetic syzygy
criterion, though infinite rings may give infinitely many such
tests. When $C$ is flat, all positive Tor groups involving
$\mathcal F$ vanish, so $K_1=0$, $K_2=W_2$, $T_1=W_1$ and the
criterion reduces to B.080, with its connecting map fixed by the
following sign comparison. Under the usual double-resolution
augmentation and the connecting convention of (B80.9),
$\delta_0=-\partial$. One may check the sign using a flat resolution
of $E$ whose degree-zero module is $\mathcal F$ and whose positive
degrees resolve $A$: the component $w_2$ of $w$ in
$Q_1\otimes\mathcal F$ satisfies $dw_2=-H$. The usual connecting
map $\tau$ therefore takes its image in $Q_1\otimes E$ to the
negative of the class used in (B80.9). Thus the vanishing tests
are identical, and negating the last cokernel identification gives
exactly (B80.2). The correction terms in (B82.19) are essential
in the nonflat scope; they are not discarded.

Here are two checks which also show how the invariants operate
without assuming flatness. For the forest inclusion of B.067,
write $D=S/R$ and use its monic basis
$b_i=X^i$ for $i<4$, $b_i=(X^4+X^2)X^{i-4}$ for $i\ge4$.
Then the coefficient and value modules in (B82.1) are
\[
 \mathcal F=\bigoplus_{i\ge0}S b_i,\qquad
 E=\bigoplus_{i<4}D b_i,
                                                        \tag{B82.20}
\]
and $\varepsilon$ is the quotient $S\to D$ on the four low
summands and zero on the other summands. This follows directly
from the Hermite value calculation in B.067. Put
$M_q=\operatorname{Tor}_q^R(S,S)$ for $q\ge1$.
The exact sequence $0\to R\to S\to D\to0$ and
$S\otimes_RS=S$ give natural isomorphisms
\[
 M_q\simeq\operatorname{Tor}_q^R(S,D)
     \simeq\operatorname{Tor}_q^R(D,S)\quad(q\ge1).
                                                        \tag{B82.21}
\]
Indeed the surrounding positive Tor groups with the free module
$R$ vanish, and in degree one the remaining map
$S\otimes_RR\to S\otimes_RS$ is an isomorphism. The same
exact sequence gives an injection
\[
 0\longrightarrow M_1\longrightarrow
 \operatorname{Tor}_1^R(D,D)\longrightarrow D\longrightarrow0
                                                        \tag{B82.22}
\]
and isomorphisms $M_q\simeq\operatorname{Tor}_q^R(D,D)$
for $q\ge2$. Tor commutes with these direct sums, so the rows
in (B82.3) decompose by pairs of basis indices. For two low
indices, $a_1$ is the diagonal of $M_1$, while $b_1$ is the
difference followed by the injection in (B82.22); hence the
row homology $K_1$ is zero. For exactly one low index, $a_1$
is an isomorphism and $W_1=0$, again giving zero row homology.
For two high indices both $V_1,W_1$ are zero. Also $b_2$ is
surjective on every pair with nonzero $W_2$, so $T_2=K_2=0$.
Consequently, in this explicitly nonflat case,
\[
                  K_1=K_2=0,\qquad T_1\simeq D^{16}.
                                                        \tag{B82.23}
\]
The connecting map records, up to its fixed minus sign, the
sixteen low-low coefficients modulo $R$. This follows either
by the lift (B82.12) and (B82.22), or by taking the coefficient
cokernel of the displayed split summands. The identities
\[
\begin{aligned}
 Q(X+Y)&=Q(X)+Q(Y),\\
 Q(XY)&=Q(X)Q(Y)+Q(X)Y^2+X^2Q(Y),\qquad Q=X^4+X^2,
\end{aligned}
\]
show that the two arrays for every $QG\in Q S[X]$ have zero
low-low remainder; those for $R[X]_{<4}$ have that remainder
over $R$. Thus both obstruction classes in (B82.7) vanish.
This verifies the positive case directly in the new invariants,
without falsely imposing $M_q=0$.

For the flat cusp of B.079, let $R=k+tV$, $V=l[[t]]$, $K=l((t))$,
$k=\mathbb F_q$, $l$ infinite, and $P=X^q-X$. Its value module
is specified explicitly by the monic basis:
\[
 E=\bigoplus_{i<q}(K/R)X^i\ \oplus\!
   \bigoplus_{\substack{i<q\\j\ge1}}
        (K/t^{1-j}V)X^iP^j.                              \tag{B82.24}
\]
In particular $K/V$ is a direct summand. Choose
$\beta\in l\setminus k$. The relation
$\beta\otimes1-1\otimes\beta$ is nonzero in
$V\otimes_RV$, because reduction modulo $tV$ sends it to
the nonzero element of $l\otimes_k l$ detected by a
$k$-linear functional taking $1$ to zero and $\beta$ to one.
It maps to zero in $K\otimes_RV=K$. The two long exact
sequences for $0\to V\to K\to K/V\to0$, with $K$ flat,
therefore give a unique nonzero class
\[
       \eta\in\operatorname{Tor}_2^R(K/V,K/V)
\]
whose double connecting image is that relation. The split
inclusion of $K/V$ in (B82.24) makes its image nonzero in
$W_2=\operatorname{Tor}_2^R(E,E)$. Since $C=K$ is flat,
this is a nonzero class in $K_2$. The criterion rejects the
cusp at the uniqueness test. Its nonzero kernel class corresponds,
with the connecting sign fixed as above, to
$\beta P\otimes P-P\otimes\beta P$ from B.079.
These calculations check both a nonflat positive inclusion and
a flat negative inclusion within the universal criterion.
\end{proof}
